% ============================================================
% introduction.tex
% ============================================================

Quantum mechanics is the most precisely verified theory in the history
of science.
It answers, with extraordinary accuracy, every question about how particles
behave: what energies they occupy, how they scatter, how they interfere,
how they decay.
But it is structurally silent on a different class of questions --- why
its own axioms hold.

Why are atomic spectra discrete rather than continuous?
The Schr\"odinger equation produces discrete eigenvalues given the right
boundary conditions, but the boundary conditions are loaded into the
formalism before the calculation begins.
Why is the Born rule $|\psi|^2 = \text{probability}$ true?
It is a postulate --- von Neumann's measurement axiom --- not a derivation.
Why do the Dirac gamma matrices have the algebraic structure they do?
Dirac constructed them to satisfy his requirements; they were
not derived from anything geometrically prior.
Why is the gauge group of the Standard Model $\mathrm{SU}(3)\times
\mathrm{SU}(2)\times\mathrm{U}(1)$ and not some other group?
The Standard Model takes its gauge content as given.

These are not criticisms of quantum mechanics.
They are a precise statement of what it is silent about, and that silence
has been known since at least von Neumann~\cite{vonneumann1932} and is
uncontroversial.
A framework that answered these questions would not replace quantum mechanics;
it would explain it.

\medskip

This paper proposes such a framework.
We show that all four questions have geometric answers, and that all four
answers follow from a single conservation law:

\begin{center}
\textbf{Conservation of probability --- the axiom $\mathcal{A}=1$.}
\end{center}

On a bipartite three-dimensional causal lattice $\mathcal{T}_\diamond^3$,
each particle is a \emph{causal session}: an independent phase rotor
carrying amplitude $(\psi_R, \psi_L)$ on the two chirally-opposite
sublattices, subject to $|\psi_R|^2 + |\psi_L|^2 = 1$ at every tick.
No Hamiltonian, no Hilbert space, and no measurement postulate are
assumed.

The four answers:
\begin{enumerate}
\item \textbf{Quantization is an attractor, not a boundary condition.}
    The discrete mass spectrum arises from Arnold tongue frequency-locking
    between the particle's internal Zitterbewegung frequency and the vacuum
    carrier.
    Any nearby frequency is dynamically pulled into the nearest rational
    resonance basin.
    Discreteness is self-organising, not postulated.

\item \textbf{The Born rule follows from $\mathcal{A}=1$.}
    Probability is amplitude squared because amplitude is the only
    physical quantity on the lattice and $\mathcal{A}=1$ is the
    conservation law that governs it.
    The Born rule is not a measurement axiom; it is a conservation theorem.

\item \textbf{The Dirac gamma matrices encode lattice geometry.}
    The bipartite RGB/CMY sublattice structure \emph{is} the two-component
    Dirac spinor.
    In the continuum limit $a\to 0$, the tick rule converges to the
    Dirac equation $(i\gamma^\mu\partial_\mu - m)\psi = 0$, with the
    gamma matrices identified as the geometric encoding of three pairs
    of antiparallel basis vectors.
    The Clifford algebra $\{\gamma^\mu,\gamma^\nu\} = 2g^{\mu\nu}$ holds
    because the basis vectors form a tight frame --- a geometric theorem,
    not an algebraic postulate.

\item \textbf{Photon emission is a conservation-law necessity.}
    The Coulomb interaction is a probability \emph{source} --- it injects
    amplitude into the electron session from outside.
    $\mathcal{A}=1$ cannot accommodate this within a single session;
    it mandates creation of a new photon session at orbital lock-in.
    Emission, recoil, and energy conservation all follow from
    $\mathcal{A}=1$ alone.

\item \textbf{Gravity is refraction, not curvature.}
    Every causal session carries an independent tick counter.
    The local density of these clocks --- the session density weighted by
    probability --- is the gravitational field.
    Causal paths refract toward regions of higher clock density by the
    same mechanism as Fermat's principle in optics: amplitude hops
    preferentially toward nodes where the phase cost is lower.
    No curved geometry is required.
    Gravitational time dilation is scheduler load: a clock near a mass
    processes more causal events per external tick and therefore runs
    slow.
    The equivalence principle is structural: both gravitational and
    inertial mass arise from the same phase-mismatch parameter $\omega$,
    so their equality requires no explanation.
    The event horizon is scheduler saturation --- a computational
    deadlock at the Planck density ceiling $\rho_\text{max} =
    \ell_P^{-3}$ --- with no singularity, and Hawking radiation as
    stochastic queue-drop.
    In the continuum limit the clock fluid equations reduce to the
    Einstein field equations; at finite lattice spacing they carry
    discrete corrections to the $1/r^2$ law that are computable and
    falsifiable.
\end{enumerate}

\medskip

A further result, which we consider the deepest, concerns the nature of
bound states.
The standard treatment of hydrogen places the proton as a fixed external
potential $V(r) = -e^2/r$ and solves for the electron's eigenstates.
This is not merely an approximation --- it is the wrong physical picture.

A static Coulomb potential is insufficient to produce stable quantised
orbits from the lattice dynamics alone.
The proton must be represented as an active session --- a phase rotor
participating in the joint $\mathcal{A}=1$ accounting every tick ---
for the electron to explore the well and lock onto resonant orbits
spontaneously.
Hydrogen is not an electron in a field.
It is \emph{two sessions finding a joint Arnold tongue attractor.}

This distinction has directly observable consequences.
In a static well the electron collapses to the potential minimum regardless
of initial momentum; the basin of attraction around the Bohr orbit has
zero width (\texttt{exp\_11} $n=2$: flat landscape across all
$k$-values tested).
With an active proton session, the proton's Zitterbewegung continuously
breaks the spherical symmetry of the potential, providing the stochastic
perturbation that drives the electron into the resonance basin;
the Bohr orbital wavevector is recovered to four significant figures
with no parameter tuning (\texttt{exp\_12}).
The basin of attraction is finite and self-correcting --- the orbit is
stable not because we initialised it exactly, but because the two-session
dynamics continuously restores it.

Photon emission then follows as a three-session event: the joint
proton-electron resonance transitions to a lower Arnold tongue, the
displaced amplitude is absorbed by a new photon session, and the proton
adjusts to the new joint configuration.
None of this requires a separately postulated transition rate.

\medskip

\paragraph{Relation to prior discrete spacetime programmes.}
The $\Tdiamond$ framework belongs to a tradition of discrete spacetime
proposals that includes causal set theory~\cite{bombelli1987},
'T~Hooft's cellular automaton interpretation~\cite{thooft2016},
Wolfram's computational universe~\cite{wolfram2002}, and Zuse's
Rechnender Raum~\cite{zuse1969}.
It differs from these in the precision of its structural claims and the
number of falsifiable numerical predictions it makes without free
parameters.
The frame condition $\sum_\text{RGB}\mathbf{v}_i\mathbf{v}_j^T = 3\delta_{ij}$
is not a tuning choice; it is a geometric identity of the RGB basis vectors
that forces the Clifford algebra and the isotropic dispersion relation
simultaneously.
The lattice does not approximate known physics --- it derives it from
a single conservation law and a single geometric substrate.

\paragraph{Organisation of the paper.}
Section~\ref{sec:substrate} defines the $\Tdiamond$ lattice, its
bipartite structure, and the calibration of physical units.
Section~\ref{sec:causal_sessions} introduces causal sessions and the
phase rotor.
Sections~\ref{sec:tick_scheduler}--\ref{sec:kinematics} derive
momentum, force, and the Dirac equation from the tick rule.
Sections~\ref{sec:gravity}--\ref{sec:vacuum_twist} treat gravity
and electromagnetism as distinct phase geometries.
Sections~\ref{sec:interference}--\ref{sec:observer} cover interference,
decoherence, and the observer.
Sections~\ref{sec:predictions}--\ref{sec:hydrogen} present the
falsifiable predictions and the hydrogen results.
Section~\ref{sec:conclusion} states what has been established and
what remains open.
